\section{INTRODUCTION}

The dynamics of rolling bodies on an incline form an essential part of the undergraduate physics syllabus. Despite being a well-known problem, even slight changes to the symmetry of the rolling body or the surface on which it is rolling can lead to interesting and unexpected dynamics, which are the subject of contemporary research. Prominent examples are provided by biased balls exhibiting various dynamic behaviors, including slow precession and nutation~\cite{Wheathead:AJP2020,Kilin:RCD2017}, surprising behavior of spherical magnets on an aluminum incline \cite{Cross:PE2025}, and dynamics of particle ensembles in viscous media \cite{Wolfschwenger:AMM2024}. The dynamics of objects with multiple rotational axes are particularly interesting. For instance, a biased ball that spins around its symmetry axis while also rotating around a vertical axis through its center of mass (COM) can exhibit complicated, curved trajectories of various shapes \cite{Jalali:PRE2015}. In most publications, both the rotational axes and their corresponding torques were considered to be of a mechanical nature. However, if a rolling body is magnetic, an additional, time-dependent magnetic torque and a corresponding axis appear in an external magnetic field. Consequently, even a perfectly shaped magnetic solid sphere could exhibit non-trivial rolling dynamics.

The first indication of this behavior was recently provided in  \cite{Vedmedenko:SciRep2022}. It has been demonstrated that if a magnetic solid sphere rolls due to some force and the magnetic and mechanical rolling torques are not collinear, the solid sphere's rolling trajectory will have a transversal component, causing it to start revolving and performing a type of planetary motion known as revolution. Although it has been pointed out in \cite{Vedmedenko:SciRep2022} that this revolutional motion depends on the relative orientation of the external magnetic field and the mechanical force, the shape of the trajectories remained unknown. 

The present investigation aims to provide a comprehensive description of the rolling dynamics of a ferromagnetic solid sphere on an incline. Our results demonstrate that the trajectory of such a sphere deviates significantly from linear rolling and exhibits revolution, with the center of the orbit lying outside the ball. However, the shape of the trajectory strongly depends on the initial mutual orientation of the magnetic moment $\hat{\bm{\mu}}$ and the external magnetic field $\bm{B}$. Interestingly, the trajectories change under exchange of $\hat{\bm{\mu}}$ and $\hat{\bm{B}}$. 
%It is neither symmetric with respect to the $\hat{\bm{\mu}}|_{t=0\unit{\second}}\times \hat{\bm{B}}$ nor with respect to the cross-product of the initial velocity and magnetic field $\hat{\bm{v}}|_{t=0 \unit{\second}} \times \hat{\bm{B}}$. 
The dynamics of the rolling motion are complex and include spinning, precession, and nutation of the magnetic moment.

The manuscript is divided into the following sections. Section 2 is devoted to the analytical derivation of the equations of motion and a brief description of the numerical procedure, while chapter 3 comprehensively analyses the phase space of the rolling trajectories. Finally, we conclude by summarizing our findings.